% Pregunta 345 (2020-enero-OEP2019, n.º 12): Shapiro-Stiglitz, condición de estímulo al esfuerzo (CEE).
% Con b > 0 la CEE es creciente y vertical al acercarse a N (en azul, de referencia).
% Con b = 0 es horizontal en w = w̄ + e + eρ/q (no depende de L) y baja si sube q.
% Autor: Víctor Gutiérrez Marcos
\documentclass[tikz,border=6pt]{standalone}
\input{tikz-tcee}
\begin{document}
\begin{tikzpicture}[tcee]
  \draw[eje] (0,0) -- (6.4,0) node[right] {$L$};
  \draw[eje] (0,0) -- (0,5.4) node[above] {$w$};
  \node[below left] at (0,0) {$0$};
  % Pleno empleo
  \draw[guia] (5.2,0) -- (5.2,5.2);
  \node[below] at (5.2,0) {$N$};
  % Caso general b > 0: creciente y vertical cerca de N
  \draw[curva2,domain=0:4.67,samples=60,opacity=0.55] plot (\x,{3.1+0.25*\x/(5.2-\x)});
  \node[etiqueta,left,text=tceeazul,align=right] at (4.3,4.9) {CEE con $b>0$\\(creciente)};
  % Caso del enunciado b = 0: horizontal
  \draw[curva1] (0,2.4) -- (5.2,2.4) node[etiqueta,above left,text=tceemorado] {CEE con $b=0$};
  \draw (0,2.4) -- (-0.08,2.4) node[left] {$\bar w+e+\dfrac{e\rho}{q}$};
  % Mayor probabilidad de detección
  \draw[curvanueva,desplazada] (0,1.2) -- (5.2,1.2) node[etiqueta,below left,text=tceerojo] {$q'>q$};
  \draw (0,1.2) -- (-0.08,1.2) node[left] {$\bar w+e+\dfrac{e\rho}{q'}$};
  \draw[flecha] (2.0,2.3) -- (2.0,1.3);
  \draw[flecha] (3.6,2.3) -- (3.6,1.3);
\end{tikzpicture}
\end{document}
